\section{Exact matching after a long common prefix}
\label{sec:matching}

\subsection{A remaining expensive local query}

The frozen repository already contains a complete polynomial algorithm for
longest common substrings of binary straight-line programs (SLPs). General
polynomial compressed LCS algorithms are also established in the
literature~\cite{compressed-lcs-classical}. The repository's
cut-pair search combines arithmetic progressions of suffix--prefix overlaps,
exact compressed equality, and periodic extension. Subsequent refinements
use the shared signed alphabet and the shared directed adjacent-pair set
to bound possible match lengths. A checked witness attaining such a bound
finishes the query immediately. These are substantial existing components;
the present change adds a specialized exact branch to that implementation.

The relevant baseline implementation is
\path{fast/fastunknot/compressed_lcs.py}. Its proofs are in
\path{synthesis/compressed_lcs.tex}, \path{synthesis/lcs_bounds.tex},
and \path{synthesis/lcs_transitions.tex}, relative to the frozen
\path{Topology/UnknotRecognition} directory. In particular, the last source
records an unresolved family with identical alphabets and identical
adjacent-pair sets. Put
\[
 P_N=(ab)^N,\qquad T=\mathtt{aabbccacbabc},\qquad
 X_N=P_NcT,\quad Y_N=P_NaT.
\]
Every ordered pair over $\{a,b,c\}$ occurs inside $T$. Consequently neither
the alphabet bound nor the directed-pair bound separates these words.
Their longest common prefix already has length $2N$, but a generic exact
search still has to prove that no longer substring occurs at a shifted
position. The baseline exhausts the two-million-work allowance on the
specified binary-power grammar when $N=2^{500}$.

The information that makes this query easier is the small amount of word
remaining after the checked common prefix: thirteen letters on each side.
An improvement on a nearly full match can start in only a short initial
interval. Rather than add another fixed-length local-factor bound, we
reduce all possible improvements to short shifts of the same known common
prefix. A finite cycle automaton checks those shifts on the SLP itself.
The resulting algorithm applies to an arbitrary common prefix, with no
assumption that it is alternating, uniform, or globally periodic.

\subsection{The diagonal reduction}

All intervals are zero-based and half-open. For finite words $X,Y$, let
$\LCS(X,Y)$ denote the length of their longest common contiguous substring,
and let $\LCP(X,Y)$ denote their longest common-prefix length. We write
$\operatorname{LCSuf}$ for the longest common-suffix length.

Suppose the exact common-prefix length is $K$, and write
\[
 |X|=K+a,\qquad |Y|=K+b,\qquad
 a,b\ge1,\qquad S=a+b-1.
\]
The case in which either word ends at position $K$ is already solved by
the prefix witness. In the new branch we assume
\begin{equation}
 K\ge S-1.
 \label{eq:near-prefix-domain}
\end{equation}
The implementation additionally requires $S\le32$. The latter is an
implementation gate, not an assumption needed for the mathematical reduction.

\begin{lemma}[A bounded set of cut alignments]
\label{lem:near-prefix-diagonals}
Let $c=a-1$. For each integer
\[
 1-a\le d\le b-1
\]
put $c_d=c+d$ and define
\[
 M_d=
 \operatorname{LCSuf}(X[0:c],Y[0:c_d])
 +\LCP(X[c:|X|],Y[c_d:|Y|]).
\]
Under \eqref{eq:near-prefix-domain},
\begin{equation}
 \LCS(X,Y)=\max\bigl(K,\max_{1-a\le d\le b-1}M_d\bigr).
 \label{eq:near-prefix-max}
\end{equation}
There are exactly $S$ values of $d$, and both cuts lie in the first
$S-1$ letters of the common prefix.
\end{lemma}

\begin{proof}
Each $M_d$ is the length of an actual common substring: concatenate the
matching suffix ending at the two cuts with the matching prefix beginning
there. Thus the right-hand side cannot exceed $\LCS(X,Y)$.

Conversely, suppose a common substring has length $m>K$ and starts at
positions $i$ in $X$ and $j$ in $Y$. Its starts satisfy
\[
 0\le i\le |X|-m\le a-1=c,
 \qquad 0\le j\le b-1.
\]
Its alignment $d=j-i$ therefore belongs to the displayed range. The cut
$c$ lies in the closed interval $[i,i+m]$: the left inequality was just
proved, and $i+m>K\ge c$. The corresponding cut $c_d$ lies at the same
position relative to the second occurrence because
$c_d-j=c-i$. Extending the two occurrences left and right from their cuts
therefore gives $M_d\ge m$.

Finally, $0\le c\le S-1$ and $0\le c_d\le S-1$. The domain assumption
places all these cuts inside the shared prefix. Counting the integer
alignments gives $(b-1)-(1-a)+1=a+b-1=S$.
\end{proof}

The alignment count is a sum of the two deficits, rather than their
product. Enumerating all possible starting pairs separately would use
$ab$ candidates. An arbitrary pair of compressed words can still have an
exponentially large $S$; the new branch is used only when its bound is
small. A failed gate invokes the existing complete algorithm without
enumerating any of these alignments.

Lemma~\ref{lem:near-prefix-diagonals} alone is insufficient for practical
speed. Evaluating the $M_d$ by general compressed LCP and LCSuf queries
can compare very long, differently aligned parses repeatedly. The
benchmark includes precisely this generic-LCE ablation. To obtain the
improvement, the shared-prefix structure must also be used inside each
alignment query.

\subsection{Short shifts are exact periodicity tests}

Let $P=X[0:K]=Y[0:K]$. The left part of $M_d$ uses at most $S-1$ letters
on each side, so it can be computed from one explicitly read short prefix
of $P$. The right part is more interesting. When $d=0$, its value is
$K-c$: the definition of the exact LCP says the next two letters differ.
Together with its left extension, this gives $M_0=K$.

For $d\ne0$, set
\[
 u=\min(c,c_d),\quad v=\max(c,c_d),\quad q=v-u=|d|,
 \quad A=P[u:v],\quad Z=P[u:K].
\]
Here $1\le q\le S-1$, and the seed $A$ is a short, explicit word. Let
$A^\infty$ denote the one-sided infinite periodic word beginning with $A$.
It is only a mathematical reference word; no infinite or long repeated
word is allocated. Define
\[
 E=\LCP(Z,A^\infty).
\]
Since $A$ is the first $q$ letters of $Z$, we have $E\ge q$.

\begin{lemma}[Exact self-comparison from one periodic prefix]
\label{lem:near-prefix-periodicity}
With the preceding notation,
\[
 \LCP(P[u:K],P[v:K])=E-q.
\]
If $E<K-u$, this is also the exact right extension in the original words
from cuts $c,c_d$. If $E=K-u$, only at most $\max(a,b)\le S$ additional
letters must be compared to obtain that right extension.
\end{lemma}

\begin{proof}
For every $0\le t<E-q$, both $Z[t]$ and $Z[t+q]$ are inside the prefix
that agrees with $A^\infty$, and their phases modulo $q$ are equal.
Hence the two substrings agree for $E-q$ letters. If $E<|Z|$, then
$Z[E]$ differs from the periodic letter at that phase, whereas
$Z[E-q]$ still equals it. This is the first mismatch in the shifted
comparison. If $E=|Z|$, the shorter substring has ended, so its full
length $|Z|-q$ is the answer.

In the first case the mismatch occurs before either relevant occurrence
leaves the known prefix $P$, so the same mismatch occurs in $X,Y$. In
the second case the larger cut coordinate has reached position $K$,
and the smaller has reached $K-q$. The word whose larger coordinate
has just reached $K$ has only its original suffix of length $a$ or $b$
remaining. Comparing the rest therefore takes at most $\max(a,b)$
letter comparisons, regardless of $K$.
\end{proof}

This lemma does not assume that $P$ has period $q$. A short violation
of that period is an exact reason to stop the shifted match. Nor does
it infer equality from a matching seed: the whole required periodic
extent is computed on the grammar. For arbitrary nonperiodic prefixes,
the state computation often finds an early violation; its correctness
does not depend on finding one quickly.

\subsection{Evaluating the cycle automaton on an SLP}

For a fixed seed $A$ of length $q$, a phase is an integer in
$\{0,\ldots,q-1\}$. For a grammar node $V$ and phase $r$, define
\[
 F(V,r)=\LCP\bigl(W_V, A[r]A[r+1]\cdots\bigr),
\]
where pattern indices are reduced modulo $q$ and $W_V$ is the word
spelled by $V$. A terminal $t$ satisfies
\[
 F(t,r)=\begin{cases}1,&t=A[r],\\0,&t\ne A[r].\end{cases}
\]
For a binary production $V=UV'$,
\begin{equation}
 F(V,r)=
 \begin{cases}
 F(U,r),&F(U,r)<|U|,\\[2pt]
 |U|+F\bigl(V',(r+|U|)\bmod q\bigr),&F(U,r)=|U|.
 \end{cases}
 \label{eq:cycle-automaton-prefix}
\end{equation}

\begin{proposition}[Exact periodic-prefix primitive]
\label{prop:cycle-automaton}
Recurrence \eqref{eq:cycle-automaton-prefix} computes the exact periodic
prefix extent. For a grammar with $g$ nodes it has at most $gq$ states.
Memoized evaluation needs $O(gq)$ arithmetic and dictionary operations,
and $O(gq)$ stored states, without expanding any nonterminal word.
\end{proposition}

\begin{proof}
The terminal rule is immediate. At a concatenation, a mismatch inside
the left child is the first mismatch of the concatenation. If the whole
left child matches, the right child starts after exactly $|U|$ periodic
letters, so its phase is $(r+|U|)\bmod q$. This proves the recurrence
by induction on the grammar DAG.

There are only $q$ possible phases for each node. The evaluator memoizes
each requested pair $(V,r)$ once. Each state requests at most its two
child states, with a constant amount of additional arithmetic and table
work. Children precede their parent in the binary SLP, so dependencies
are acyclic. Repeated occurrences in an exponentially large derivation
tree reuse the same node--phase entry; they do not produce new states.
\end{proof}

This is an elementary compressed finite-automaton computation, used here
to replace a more expensive general equality primitive. We do not claim
that evaluating finite automata on compressed words is a new general
idea. The present contribution is the exact alignment reduction that
produces small cycle automata for this particular LCS regime, together
with its implementation, proof, and controlled measurements.

\subsection{The resulting algorithm and bit complexity}

\begin{theorem}[LCS with bounded suffix deficits]
\label{thm:near-prefix-lcs}
Suppose binary SLPs for $X,Y$ and the exact value $K=\LCP(X,Y)$ are
available, with positive suffix deficits $a,b$, $S=a+b-1$, and
$K\ge S-1$. Let $g$ be the number of allocated nonempty grammar nodes
at entry to the algorithm, and let $H\le g$ be their maximum height.
There is an exact algorithm returning the LCS length and an occurrence
witness using
\[
 O(gS^2)
\]
arithmetic and dictionary operations. It creates $O(SH)$ additional
grammar nodes, has at most $O(gS)$ simultaneously live node--phase
states, and explicitly reads at most $5S-1$ letters from the roots.

For fixed $S$, the implemented branch has deterministic polynomial bit
complexity, in addition to the preceding exact LCP computation. It does
not require a random fingerprint, a bound on hash-collision probability,
or a promise of periodicity of the common prefix.
\end{theorem}

\begin{proof}
Read $P[0:S-1]$ and the suffix windows
\[
 X[\max(0,K-S):|X|],\qquad
 Y[\max(0,K-S):|Y|].
\]
Each suffix window has length at most $2S$, giving the stated total
letter bound. The routine follows only grammar branches intersecting
these bounded intervals; a conservative traversal bound is $O(HS)$.

Enumerate the $S$ alignments from Lemma~\ref{lem:near-prefix-diagonals}.
Their left extensions are computed from the short prefix already read.
The zero shift needs no periodic query. For every nonzero shift construct
the slice $P[u:K]$, take its seed from the short prefix, and apply
Proposition~\ref{prop:cycle-automaton}. A slice of one original root adds
$O(H)$ boundary-path nodes and does not increase height. Each individual
slice has a reachable grammar of size $O(g+H)=O(g)$, even though slices
from different alignments are retained in the arena. Consequently the
sum of node--phase work is bounded by
\[
 O\!\left(g\sum_{d=1-a}^{b-1}\max(1,|d|)\right)=O(gS^2).
\]
The $O(S^2)$ short comparisons and $O(SH)$ slice work fit this bound.
At any moment only one periodic-query table is live, with at most
$O(gS)$ states. Accumulated slice allocation is $O(SH)$. When a periodic
match reaches the prefix endpoint, Lemma~\ref{lem:near-prefix-periodicity}
shows that the required continuation lies in the two already read short
suffix windows. Equation~\eqref{eq:near-prefix-max} proves that the maximum
of the resulting exact extensions and the original prefix witness is the
exact LCS. Its two starting positions are the two cuts minus the common
left-extension length.

For bit accounting, let $B$ bound the binary lengths of stored word
lengths and coordinates, let $A_0$ bound the bit length of an alphabet
symbol, and let $R$ bound the bit lengths of the caller's numeric caps
and the preexisting resource counters. Put
\[
 \beta=B+A_0+R+\lceil\log_2(gS+2)\rceil+1.
\]
All stored lengths use $O(\beta)$ bits. Additions, comparisons, table-key
operations, resource-accounting updates, and phase remainders have
polynomial bit cost; a conservative
schoolbook bound is $O(\beta^2)$ per arithmetic operation. If balanced
comparison dictionaries are used, the same state recurrence therefore
has bit bound
\[
 O\bigl(gS^2\beta^2\log(gS+2)\bigr).
\]
The delivered Python code instead uses ordinary collision-resolving
dictionaries. No probabilistic assumption is needed for exactness. Even
charging the worst-case linear lookup cost in a dictionary with
$O(gS)$ entries gives the conservative deterministic bound
\[
 O(g^2S^3\beta^2).
\]
This also covers the existing exact rule-interning and slice dictionaries
under the same conservative charge. Stored grammar data, short words,
and simultaneously live automaton states use $O(gS\beta)$ bits, apart
from implementation object overhead. With the fixed gate $S\le32$,
all these bounds are polynomial in the encoded input size. The common
prefix itself is obtained by the repository's previous exact SLP LCP
routine; its cost must be added, not silently treated as free.
\end{proof}

The two arithmetic bounds should not be confused. The useful operation
count $O(gS^2)$ assumes ordinary constant-time dictionary operations.
The explicitly stated deterministic bit bound remains polynomial even
under worst-case dictionary behavior. Both bounds concern this one
string query after its common prefix is available; neither bounds the
number of group transformations in a full recognition attempt.

\subsection{Integration, caps, and resource semantics}

The implementation adds \code{small_interval} and
\code{periodic_prefix} to \code{CommonSubstring}, with
\code{near_prefix} coordinating the branch.
The public \code{longest} method still begins with its exact prefix
query and its alphabet and directed-pair bounds. If those bounds have
not already completed the query, it calls the new branch with the exact
prefix length. An eligible branch returns an exact answer and witness;
an ineligible branch returns no answer and the existing cut-pair search
continues.

The current gate is deliberately fixed at $S\le32$. It therefore
adds no loop proportional to an exponentially large encoded deficit.
Its second condition, $K\ge S-1$, ensures that every small seed and
every selected cut lies in the known common prefix. Inputs outside this
domain are covered by the existing complete matcher. Increasing the
gate is a future tuning decision supported by the parameterized bound;
it is not required for correctness of the present integration.

An optional LCS cap is handled exactly as before. If the prefix already
attains it, the earlier code returns immediately. Otherwise the computed
prefix is the true uncapped LCP, because its mismatch occurred strictly
before the cap. The new branch may stop after an actual witness attains
the cap; it returns the first capped part of that witness. It never
uses a guessed upper bound as evidence that a substring exists.

Every traversal and comparison in the added branch consumes the arena's
shared work allowance and calls its cancellation callback. Allocation
continues to use the existing node limit. Each periodic-query memo table
also has a state allowance equal to \code{max_nodes}; exceeding it raises
\code{CompressedLimit}. These exceptions remain inconclusive resource
outcomes. They are not converted into a negative substring answer or a
negative knot verdict. The periodic table is local to one completed
query, so an interruption stores no partially proved answer for reuse.

The new short-interval reader explicitly materializes bounded seeds and
tail windows. Calling this ``no expansion'' without qualification would
hide that detail. The precise guarantee is that at most $5S-1$ root
letters are read explicitly, irrespective of the expanded core length.
No diagnostic whole-word expansion, exponential repeated-word allocation,
or probabilistic equality shortcut is used.

\subsection{Resolved families and a retained negative control}

\begin{example}[Equal adjacent-pair sets]
For $X_N=P_NcT$ and $Y_N=P_NaT$, with $N\ge16$, the exact LCP is
$K=2N$, and $a=b=13$, so $S=25$. We have
\[
 \LCS(X_N,Y_N)=2N.
\]
Indeed, a longer substring would start in the first thirteen positions
of each word. Starts of opposite parity mismatch immediately in the
alternating prefixes. For equal parities and unequal starts, the earlier
arrival at $X_N$'s tail encounters $c$ opposite an alternating letter.
If $Y_N$ reaches its tail first, its initial $a$ may match, but the next
letter is the first $a$ of $T$ and mismatches the alternating prefix's
next $b$. The equal-parity separation is at least two, so the other
word has not left that prefix before this mismatch. Equal starts
encounter $c$ versus $a$ at position $2N$. Thus no alignment permits a
longer match.

The new branch verifies this answer without occurrence tables. On the
standard power grammars its arithmetic cost is linear in
$h=\log_2N$: $S$ is fixed, and the grammar has $O(h)$ nodes. Its bit
cost remains polynomial in $h$, with the factors for arithmetic on
$O(h)$-bit lengths included as in Theorem~\ref{thm:near-prefix-lcs}.
\end{example}

\begin{example}[A genuine improvement on the prefix witness]
Let
\[
 X=(ab)^Nc,\qquad Y=(ab)^{N+1}c.
\]
Their exact LCP is $2N$, but $X$ occurs in $Y$ at offset two, so
\[
 \LCS(X,Y)=2N+1.
\]
Here the two suffix deficits are one and three. The new branch returns
the witness $(2N+1,0,2)$. This tests discovery of a longer common
substring, rather than merely certification that an existing prefix is
optimal. Signed random-word tests also exercise longer shifted matches
and swapped input order.
\end{example}

\begin{example}[A branch that remains unresolved]
Replace $T$ in the first family by $T^3$. The tail lengths are now 37,
so $S=73>32$ and the new branch is not entered. For the tested
$N\ge256$, the same parity and first-departure argument proves that the
LCS is still $2N$. Alphabets and adjacent-pair sets still agree, so the
existing complete matcher remains responsible for this control.
Resource exhaustion on its larger cases is retained in the results.
No time ratio is assigned to an unfinished baseline or an unfinished
new query.
\end{example}

The tests compare 500 periodic-prefix queries with explicit oracles and
450 nearly full LCS instances with literal substring search, including
more than one hundred cases whose best witness is longer than the
initial prefix. They also check large powers through $2^{500}$,
cap handling, returned positions, the eligibility gate, state exhaustion,
work exhaustion, cancellation, and a successful retry after a stopped
periodic query. The existing exhaustive and randomized compressed-LCS
and cyclic-overlap tests remain in place.

The accompanying experiments compare the frozen implementation with an
identical A/A control, the exact diagonal reduction using general LCE
primitives, and the full cycle-automaton branch. They additionally retain
an interior-match control already settled by the older transition bound.
The genuine knot-corpus queries measure whole recognition and independent
certificate replay, not only this local primitive. Their unchanged
certificate digests and lack of calls to the new branch support
preservation of those existing paths; they do not demonstrate a speedup
in completed unknot recognition.

\subsection{What this result does and does not establish}

The exact theorem resolves an identifiable compressed-string bottleneck
that survived alphabet and directed-pair pruning. Its useful parameter is
the amount of input remaining after an exact long common prefix. This is
different information from a fixed-order local factor set, and the
algorithm handles arbitrary contents of the shared prefix within the
stated domain.

The result is also deliberately local. The implementation is a specialized
branch inside an already polynomial compressed LCS algorithm; it is not
a new proof that general compressed LCS belongs to polynomial time.
It does not show that all knot presentations reach this favorable regime,
that all useful cyclic overlaps have a long common prefix in the spelling
being queried, or that word grammars remain small throughout a search.
In particular, doubling a pair of relators for a cyclic query can make
their suffix deficits large even when their original spellings share
a long prefix. The new branch then yields to the existing complete
cyclic machinery. Turning the present local improvement into a bound
for general unknot recognition still requires separate global progress,
representation-size, and completeness arguments.
